(1) 在非磁性、正入射、实折射率近似下，介质 \(i\) 到 \(j\) 的电场反射振幅 \(r_{ij}=(n_i-n_j)/(n_i+n_j)\)。因此空气–膜反射为 \(r_{01}=(n_0-n_1)/(n_0+n_1)\)，膜–衬底反射为 \(r_{12}=(n_1-n_s)/(n_1+n_s)\)。符号表示反射场相位，不能只比较两个界面的强度反射率大小。
(2) 膜内一次往返传播积累相位 \(2\varphi=4\pi n_1d/\lambda\)，其中单程相位 \(\varphi=2\pi n_1d/\lambda\)。把第二界面产生的反射波及其后无限次往返叠加，得到总振幅 \(r=(r_{01}+r_{12}e^{2i\varphi})/(1+r_{01}r_{12}e^{2i\varphi})\)。这也可由双界面传输矩阵推出，分母反映多次内反射。
(3) 若 \(n_0<n_1<n_s\)，两个单界面的 \(r_{01},r_{12}\) 均为负实数。为使分子为零，需 \(r_{01}=r_{12}\) 且 \(e^{2i\varphi}=-1\)。前者化简为 \(n_1^2=n_0n_s\)，后者的最薄正膜厚为 \(d=\lambda_0/(4n_1)\)。因此常见的单层四分之一波增透膜要求折射率取几何平均，缺少任一条件都不能实现理想单色零反射。
(4) 固定膜厚时传播相位随 \(1/\lambda\) 改变，\(\lambda\ne\lambda_0\) 一般不再满足 \(e^{2i\varphi}=-1\)；材料色散又使 \(n_1^2=n_0n_s\) 不再恰好成立。如果膜有吸收，\(n_1\) 为复数，多次反射振幅和界面匹配条件同时改变，反射变小也未必意味着透射变大，因为能量可被吸收。宽带减反通常需要多层膜或梯度折射率设计。
